% ========================================================================
% CodeX-Verify: Multi-Agent Verification of LLM-Generated Code
% Complete LaTeX Document for Overleaf
% ========================================================================

\documentclass[10pt,twocolumn]{article}

% ========================================================================
% PACKAGES
% ========================================================================

% Page layout - two column for compression
\usepackage[margin=0.75in]{geometry}
\usepackage{setspace}

% Math and algorithms
\usepackage{amsmath, amssymb, amsthm}
\usepackage{algorithm}
\usepackage{algpseudocode}
\usepackage{enumerate}

% Tables and figures
\usepackage{booktabs}
\usepackage{multirow}
\usepackage{graphicx}
\usepackage{xcolor}

% References and citations
\usepackage{natbib}
\usepackage{hyperref}
\usepackage{url}

% Formatting
\usepackage{verbatim}
\usepackage{enumitem}

% Custom colors
\definecolor{ForestGreen}{RGB}{34,139,34}

% Theorem environments
\newtheorem{theorem}{Theorem}
\newtheorem{proposition}[theorem]{Proposition}
\newtheorem{lemma}[theorem]{Lemma}
\newtheorem{corollary}[theorem]{Corollary}
\newtheorem{definition}{Definition}

% Custom commands (removed \checkmark - already defined in some packages)

% Hyperref setup
\hypersetup{
    colorlinks=true,
    linkcolor=blue,
    filecolor=magenta,
    urlcolor=cyan,
    citecolor=blue,
}

% Line spacing - single space for compression
\singlespacing

% Compress spacing
\setlength{\textfloatsep}{10pt plus 2pt minus 2pt}
\setlength{\floatsep}{10pt plus 2pt minus 2pt}
\setlength{\intextsep}{10pt plus 2pt minus 2pt}
\setlength{\abovecaptionskip}{5pt}
\setlength{\belowcaptionskip}{5pt}

% Compress section spacing
\usepackage{titlesec}
\titlespacing*{\section}{0pt}{10pt}{5pt}
\titlespacing*{\subsection}{0pt}{8pt}{4pt}
\titlespacing*{\subsubsection}{0pt}{6pt}{3pt}

% ========================================================================
% TITLE AND AUTHORS
% ========================================================================

\title{\textbf{Multi-Agent Code Verification via Information Theory}}

\author{
    Shreshth Rajan \\
    Noumenon Labs, Harvard University \\
    \texttt{shreshthrajan@college.harvard.edu}
}

\date{October 2025}

% ========================================================================
% DOCUMENT START
% ========================================================================

\begin{document}

\twocolumn[
\maketitle

\begin{abstract}
LLMs generate code that passes basic tests but hides bugs: 29.6\% of SWE-bench ``solved'' patches are behaviorally wrong, and 62\% of BaxBench backend solutions are vulnerable or incorrect. Static analyzers catch these cheaply but each checks one concern at a time. We built \textsc{CodeX-Verify}, a static multi-agent verifier that runs four specialized agents (correctness, security, performance, style) in parallel in under 200ms per sample with no test execution. On 99 code samples with verified labels, the system reaches 76.1\% true-positive rate and 68.7\% accuracy. On the same samples, Semgrep with four rule packs reaches 4.8\% TPR and 28.6\% accuracy, so coordinating agents recovers most of the bugs a single static analyzer misses. An ablation over all 15 agent subsets shows the gain is concentrated in one agent: Correctness alone reaches 75.9\% accuracy, and adding Security, Performance, or Style trades a small accuracy drop for coverage of security and resource bugs the correctness checks miss. The system carries a 50\% false-positive rate because it flags quality defects (missing exception handling, weak edge-case coverage) that do not change program output; this is the central cost of static-only verification and the main limitation. On 300 Claude Sonnet 4.5 patches without ground truth it runs at 17ms per patch.

\vspace{0.2cm}
\noindent\textbf{Keywords:} Multi-agent systems, Code verification, LLM-generated code, Information theory
\end{abstract}
]


% ========================================================================
% NOTE: PASTE CONTENT FROM paper_title_abstract.tex HERE
% ========================================================================

% Starting from Section 1 (Introduction) through Appendix
% Copy from line 56 onwards in paper_title_abstract.tex


\section{Introduction}

LLMs generate code that looks correct but often fails in production. While LLM-generated code passes basic syntax checks and simple tests, recent studies show it contains hidden bugs. Xia et al.~\cite{xia2025swebench} find that 29.6\% of patches marked ``solved'' on SWE-bench don't match what human developers wrote, with 7.8\% failing full test suites despite passing initial tests. SecRepoBench reports that LLMs write secure code <25\% of the time across 318 C/C++ tasks~\cite{dilgren2025secrepobench}, and BaxBench finds 62\% of backend code has vulnerabilities or bugs~\cite{vero2025baxbench}. Studies suggest 40--60\% of LLM code contains undetected bugs~\cite{jimenez2024swebench}, making automated deployment risky.

\textbf{The Problem.} A single static analyzer checks one concern at a time and misses bugs that need a different view of the code. On our 99-sample set, Semgrep run with four rule packs (security, correctness, best-practices, Python) catches only 4.8\% of bugs at 28.6\% accuracy: it is tuned for precision and stays silent on logic errors, weak edge-case handling, and resource misuse. Test-based methods like Meta Prompt Testing~\cite{wang2024metamorphic} reach a low false-positive rate (8.6\%) by running code variants and comparing outputs, but need execution infrastructure and miss bugs that do not change output, such as SQL injection. LLM review systems like AutoReview~\cite{autoreview2025} improve security F1 by 18.72\% but check only security.

\textbf{Our Approach.} \textsc{CodeX-Verify} runs four specialized agents in parallel: Correctness (logic errors, edge cases, exception handling), Security (OWASP Top 10, CWE patterns, secrets), Performance (algorithmic complexity, resource leaks), and Style (maintainability, documentation). Each agent targets a largely disjoint bug class. Adding an informative agent never lowers the information the system has about bug presence; this monotonicity (Proposition~\ref{prop:monotone}) is the only theoretical claim we make, and whether the agents are close to conditionally independent is an assumption we state but do not measure. The design question the ablation answers is empirical: which agents actually carry the detection, and what does each added agent cost.

\textbf{Results.} We tested on 99 code samples with verified labels covering 16 bug categories mirroring SWE-bench failures. The system reaches 76.1\% TPR and 68.7\% accuracy. The head-to-head is against Semgrep on the identical samples: 4.8\% TPR, 28.6\% accuracy. Coordinating four agents recovers most of the bugs the single analyzer misses, at the cost of a 50\% false-positive rate: the system flags quality defects (missing exception handling, weak edge-case coverage) that do not change program output. This false-positive rate is the central limitation and the price of static-only verification.

The ablation over all 15 agent subsets locates the gain. Correctness alone reaches 75.9\% accuracy; Security, Performance, and Style alone stay below 21\% because each fires on a narrow bug class. Any subset without Correctness scores under 25\%. Adding Security, Performance, or Style to Correctness trades a few points of accuracy for coverage of security vulnerabilities and resource leaks that Correctness alone misses (100\% of resource bugs, 87.5\% of security bugs caught with the full system).

\textbf{Contributions.}
\begin{enumerate}[leftmargin=*, itemsep=0pt]
    \item A static multi-agent verifier that reaches 76.1\% TPR and 68.7\% accuracy on 99 verified-label samples in under 200ms per sample, with no test execution. On the same samples Semgrep reaches 4.8\% TPR and 28.6\% accuracy.

    \item An ablation over all 15 agent subsets that attributes the detection: Correctness carries the accuracy (75.9\% solo), while Security, Performance, and Style add coverage of narrow bug classes (resource, injection) that correctness checks miss.

    \item A monotonicity statement (Proposition~\ref{prop:monotone}): adding an informative agent does not reduce the system's information about bug presence. Conditional independence of agents is stated as an assumption, not measured.

    \item An open-source 99-sample dataset with verified labels, plus the full ablation and the paired Semgrep baseline on those samples.
\end{enumerate}

% ========================================================================
% SECTION 2: RELATED WORK
% ========================================================================

\section{Related Work}

\subsection{LLM Code Generation and Verification}

SWE-bench~\cite{jimenez2024swebench} evaluates LLMs on 2,294 real GitHub issues across 12 Python repositories. Follow-up work found problems: Xia et al.~\cite{xia2025swebench} show 29.6\% of ``solved'' patches don't match what developers wrote, with 7.8\% failing full test suites despite passing initial tests. OpenAI released SWE-bench Verified~\cite{openai2024verified}, a 500-sample subset with human-validated labels. Security benchmarks show worse results: SecRepoBench~\cite{dilgren2025secrepobench} reports <25\% secure code across 318 C/C++ tasks, and BaxBench~\cite{vero2025baxbench} finds 62\% of 392 backend implementations have vulnerabilities or bugs (only 38\% are both correct and secure). Across benchmarks, 40--60\% of LLM code contains bugs.

Wang and Zhu~\cite{wang2024metamorphic} propose Meta Prompt Testing: generate code variants with paraphrased prompts and detect bugs by checking if outputs differ. This achieves 75\% TPR with 8.6\% FPR on HumanEval. It requires test execution infrastructure and misses security vulnerabilities (SQL injection produces consistent outputs despite being exploitable) and quality issues that don't affect outputs. AutoReview~\cite{autoreview2025} uses three LLM agents (detector, locator, repairer) to find security bugs, improving F1 by 18.72\% on ReposVul, but only checks security, not correctness or performance. We differ by (1) coordinating four static agents that span correctness, security, performance, and style, (2) attributing the detection with an ablation over all 15 agent subsets, and (3) reaching 76.1\% TPR via static analysis without test execution.

\subsection{Multi-Agent Systems for Software Engineering}

He et al.~\cite{he2024multiagent} survey 41 LLM-based multi-agent systems for software engineering, finding agent specialization (requirement engineer, developer, tester) as a common pattern. Systems like AgentCoder, CodeSIM, and CodeCoR use multiple agents to \emph{generate} code collaboratively, but focus on producing code rather than checking it for bugs. MAGIS~\cite{magis2024} uses 4 agents to solve GitHub issues, but measures solution quality (pass@k) rather than bug detection.

Prior multi-agent work for software engineering generates code rather than verifying it, and the one security-focused reviewer (AutoReview) checks a single concern and reports no ablation. We cover correctness, security, performance, and style, and attribute each agent's contribution with an ablation over all 15 subsets.

\subsection{Static Analysis and Vulnerability Detection}

Static analysis tools (SonarQube, Semgrep, CodeQL, Checkmarx) use pattern matching and dataflow analysis to find bugs without running code. Benchmarks~\cite{johnson2024sast} show 65\% average accuracy with 35--40\% false positives, though results vary: Veracode claims <1.1\% FPR on curated enterprise code~\cite{veracode2024}, while Checkmarx shows 36.3\% FPR on OWASP Benchmark~\cite{owasp2024}. These tools check one thing at a time (security patterns, code smells, complexity) without combining analyses. Semgrep Assistant~\cite{semgrep2025} uses GPT-4 to filter false positives, reducing analyst work by 20\%, but still runs as a single agent.

Neural vulnerability detectors~\cite{ding2024vulnerability} use Graph Neural Networks and fine-tuned transformers (CodeBERT, GraphCodeBERT) trained on CVE data, achieving 70--80\% accuracy. They need large training sets (10K+ samples), inherit bias toward historical vulnerability types, and lack interpretability for security decisions. Our static analysis is deterministic and explainable without needing training data, though with higher false positives than learned models.

We extend static analysis by coordinating multiple agents that check different bug types, running them against a paired single-analyzer baseline (Semgrep) on identical samples, and attributing each agent's contribution through an ablation over all 15 subsets.

\subsection{Ensemble Learning and Information Theory}

Dietterich~\cite{dietterich2000ensemble} shows that ensembles of classifiers beat individuals when base learners are accurate and make errors on different inputs. Breiman's bagging~\cite{breiman1996bagging} and boosting~\cite{schapire1990boosting} confirm this, with theory showing ensemble error decreases as $O(1/\sqrt{n})$ for uncorrelated errors. Code verification differs from standard ensemble ML by having non-i.i.d. bug distributions, class imbalance (roughly 2.5:1 buggy:good here), and asymmetric costs (missing bugs vs. false alarms).

Cover and Thomas~\cite{cover2006information} define mutual information as $I(X; Y) = H(Y) - H(Y|X)$ and establish that conditioning on more variables does not increase conditional entropy, so adding a source never decreases the mutual information with the target. This monotonicity is the one information-theoretic fact we rely on (Proposition~\ref{prop:monotone}).

Sheyner et al.~\cite{sheyner2002attack} model multi-step network exploits as attack graphs (directed graphs of vulnerability chains). Later work~\cite{bayesian2018attack} adds Bayesian risk and CVSS scoring~\cite{cvss2024}. Attack graphs focus on network vulnerabilities (host compromise, privilege escalation) rather than code-level bug detection, which our work addresses through multi-agent static analysis.

% ========================================================================
% SECTION 3: THEORETICAL FRAMEWORK
% ========================================================================

\section{Theoretical Framework}

The framework is deliberately thin. We state the one fact that holds without further assumptions (adding an informative agent cannot reduce the system's information about bug presence), name the assumption under which that fact becomes a strict gain (conditional independence of agents), and then let the ablation in Section~6 decide empirically which agents carry the detection. We do not derive sample-size guarantees; with $n=99$ they would be vacuous.

\subsection{Problem Formulation}

Let $\mathcal{C}$ be the space of code samples and $B \in \{0, 1\}$ indicate bug presence (1 = buggy, 0 = correct). Each agent $i \in \{1, 2, 3, 4\}$ analyzes code $c \in \mathcal{C}$ through domain-specific function $\phi_i: \mathcal{C} \to \mathcal{O}_i$, producing observation $A_i = \phi_i(c)$ and decision $D_i \in \{0, 1\}$. We want aggregation function $\psi: \{D_1, D_2, D_3, D_4\} \to \{0, 1\}$ that maximizes bug detection while minimizing false alarms:
\begin{equation}
    \max_\psi \; \mathbb{P}[D_{\text{sys}} = 1 \mid B = 1] \quad \text{subject to} \quad \mathbb{P}[D_{\text{sys}} = 1 \mid B = 0] \leq \epsilon
\end{equation}
where $D_{\text{sys}} = \psi(D_1, D_2, D_3, D_4)$ and $\epsilon$ is acceptable false positive rate. This captures the tradeoff: maximize true positive rate (TPR) while keeping false positive rate (FPR) below $\epsilon$.

\subsection{Monotonicity of the System's Information}

\textbf{Idea.} An extra agent is an extra observation of the code. Conditioning on more observations never increases uncertainty about the label, so the information the system has about bug presence cannot go down when an agent is added. This is the only guarantee the architecture buys for free; it says the system is at least as informative as its best single agent, not that it is strictly better in accuracy.

\begin{proposition}[Information monotonicity]
\label{prop:monotone}
For any agent set $S$ and any agent $i$, $I(A_{S \cup \{i\}}; B) \geq I(A_S; B)$. In particular $I(A_1, \ldots, A_n; B) \geq \max_i I(A_i; B)$.
\end{proposition}

\begin{proof}
$I(A_{S \cup \{i\}}; B) - I(A_S; B) = I(A_i; B \mid A_S) \geq 0$, since conditional mutual information is non-negative~\cite{cover2006information}. Applying this along any ordering of the agents gives the second statement.
\end{proof}

The inequality is weak in general and becomes strict only when the added agent carries information about $B$ not already present in $S$, i.e., $I(A_i; B \mid A_S) > 0$. Whether the agents are close to conditionally independent given $B$, which would make each addition contribute near its full solo information, is an \emph{assumption} we state and do not test: computing $I(A_i; A_j \mid B)$ reliably needs more labeled samples than the 99 we have. The ablation in Section~6 measures the realized gains directly instead of inferring them from an independence assumption, and those gains are what the architecture claim rests on.

\subsection{Weighted Aggregation}

\textbf{Weight Selection Heuristic.} We set agent weights from two factors: individual accuracy $p_i$ and domain criticality $\gamma_i$, with $w_i \propto p_i \cdot \gamma_i$. Higher-accuracy agents receive higher weight; the criticality term $\gamma_i$ captures asymmetric costs, since a security bug blocks deployment more than a style issue.

We set $w = (0.45, 0.35, 0.15, 0.05)$ for (Security, Correctness, Performance, Style). Security gets the highest weight (0.45) despite 20.7\% solo accuracy because a missed vulnerability is costly ($\gamma_{\text{sec}}$ high). Correctness has the highest solo accuracy (75.9\%) and second weight (0.35). Performance and Style (17.2\% solo) get 0.15 and 0.05. These weights are hand-tuned on the evaluation set, not learned.

\subsection{Agent Selection}

We partition bugs into four classes: $\mathcal{B}_{\text{corr}}$ (logic errors, edge cases, exception handling), $\mathcal{B}_{\text{sec}}$ (injection, secrets, unsafe deserialization), $\mathcal{B}_{\text{perf}}$ (complexity, scalability, resource leaks), and $\mathcal{B}_{\text{style}}$ (maintainability, documentation). These classes are nearly disjoint by construction: SQL injection, an off-by-one error, and $O(n^2)$ complexity are detected by different checks. This near-disjointness is the design intuition behind the conditional-independence assumption in Proposition~\ref{prop:monotone}; we do not measure the agents' conditional correlations.

\subsection{Decision Function}

Aggregated score: $S_{\text{sys}} = \sum_{i=1}^4 w_i \cdot S_i$. Decision:
\begin{equation}
    D_{\text{sys}} = \begin{cases}
        \text{FAIL} & \text{if } |\mathcal{I}_{\text{crit}}| > 0 \\
        \text{FAIL} & \text{if } |\mathcal{I}_{\text{high}}^{\text{sec}}| \geq 1 \text{ or } |\mathcal{I}_{\text{high}}^{\text{corr}}| \geq 2 \\
        \text{WARNING} & \text{if } S_{\text{sys}} \in [0.50, 0.85] \text{ or } |\mathcal{I}_{\text{high}}| \geq 1 \\
        \text{PASS} & \text{otherwise}
    \end{cases}
\end{equation}

Security blocks on 1 HIGH, correctness on 2 HIGH, style never blocks. WARNING allows human review for borderline cases.

\subsection{What the Theory Does and Does Not Claim}

Proposition~\ref{prop:monotone} guarantees only that the full system is at least as informative about bug presence as its best single agent. It does not predict a specific accuracy gain, and it does not establish strict improvement without the (untested) conditional-independence assumption. The actual per-agent gains are an empirical result of the ablation (Section~6.2), not a theoretical prediction. We make no sample-complexity or generalization claim: at $n=99$ a uniform-convergence bound over 15 configurations is vacuous, and we report the bootstrap confidence interval ($\pm 9.1$pp) directly instead.

% ========================================================================
% SECTION 4: SYSTEM DESIGN
% ========================================================================

\section{System Design}

\subsection{Architecture}

\textsc{CodeX-Verify} runs a pipeline: code input $\to$ parallel agent execution $\to$ weighted result aggregation $\to$ deployment decision (Figure~\ref{fig:architecture}).

\begin{figure*}[t]
\centering
\small
\begin{verbatim}
                    Code Input
                        |
                AsyncOrchestrator
               /       |      |      \
        Correctness Security Perf  Style
        (AST, edge  (Vuln    (Algo  (Quality
         cases)     patterns) complex) metrics)
              \       |      |      /
                  Aggregator
              (Weighted score combination)
                        |
                  Decision Logic
            (FAIL / WARNING / PASS)
\end{verbatim}
\caption{System architecture: parallel multi-agent analysis.}
\label{fig:architecture}
\end{figure*}

Design: (1) Agents check disjoint bug classes (correctness, security, performance, style). (2) Run in parallel via \texttt{asyncio}, <200ms latency. (3) Combine weighted scores and make a deployment decision.

\subsection{Agent Specializations}

\subsubsection{Correctness Critic (Solo: 75.9\% Accuracy)}

Checks: complexity (threshold 15), nesting depth (4), exception coverage (80%), input validation (70%), resource safety (\texttt{with} statements), edge cases (\texttt{None}, empty, zero checks), contract validation (docstring vs. code), dangerous ops (\texttt{eval}, \texttt{exec}). Score: $S_{\text{corr}} = \max(0.3, 1.0 - \text{penalty}/\text{threshold})$ with severity weights (CRITICAL: 1.0, HIGH: 0.7, MEDIUM: 0.3, LOW: 0.1).

\subsubsection{Security Auditor (Solo: 20.7\% Accuracy)}

Patterns (15+, CWE/OWASP): SQL injection, command injection (\texttt{os.system}), code execution (\texttt{eval}, \texttt{exec}), unsafe deserialization (\texttt{pickle.loads}), weak crypto (\texttt{md5}, \texttt{sha1}). Secrets via regex (AWS keys, GitHub tokens, API keys, 11 patterns) and entropy ($H(s) = -\sum_i p_i \log_2 p_i$; threshold 3.5). SQL injection near \texttt{password} escalates HIGH → CRITICAL (multiplier 2.5).

\subsubsection{Performance \& Style (Solo: 17.2\% each)}

Performance checks: loop depth (0→$O(1)$, 1→$O(n)$, 2→$O(n^2)$, 3+→$O(n^3)$), recursion (tail ok, exponential flagged), anti-patterns (string concatenation in loops, nested searches). Context-aware: patch mode (<100 lines) uses 1.5$\times$ tolerance multipliers.

Style checks: Halstead complexity, naming (PEP 8), docstring coverage, comment density, import organization. All style issues LOW severity (never blocks), preventing 40\% FPR from style-based blocking.

\subsection{Aggregation}

Agents run in parallel via \texttt{asyncio} (150ms max vs. 450ms sequential). Aggregation: collect issues from all agents, adjust severities based on code context, merge duplicate detections, compute weighted score $S_{\text{sys}} = \sum_i w_i S_i$, and apply decision thresholds (Algorithm~\ref{alg:decision}).

\subsection{Decision Logic}

\begin{algorithm}[t]
\caption{Deployment Decision}
\label{alg:decision}
\begin{algorithmic}[1]
\Procedure{DecideDeployment}{$S_{\text{sys}}$, $\mathcal{I}$}
    \If{critical issues detected}
        \State \Return FAIL
    \ElsIf{security HIGH $\geq 1$ or correctness HIGH $\geq 2$}
        \State \Return FAIL
    \ElsIf{correctness HIGH $= 1$ or score $\in [0.50, 0.85]$}
        \State \Return WARNING
    \ElsIf{score $\geq 0.70$ and no HIGH issues}
        \State \Return PASS
    \Else
        \State \Return WARNING
    \EndIf
\EndProcedure
\end{algorithmic}
\end{algorithm}

Security blocks on 1 HIGH, correctness on 2 HIGH, performance/style never block alone. WARNING defers borderline cases to human review.

Calibration: initial thresholds gave 75\% TPR, 80\% FPR. Changes: (1) style MEDIUM → LOW (-40pp FPR), (2) allow 1 security HIGH (+5pp TPR), (3) weights (0.45, 0.35, 0.15, 0.05) vs. uniform (0.25 each) improved F1 from 0.65 to 0.78. Final: 76.1\% TPR, 50\% FPR.

% ========================================================================
% SECTION 5: EXPERIMENTAL EVALUATION
% ========================================================================

\section{Experimental Evaluation}

\subsection{Dataset}

We curated 99 samples: 29 hand-crafted mirroring SWE-bench failures~\cite{jimenez2024swebench, xia2025swebench} (edge cases, security, performance, resource leaks), 70 Claude-generated (90\% validation rate). Labels: buggy (71), correct (28). Categories: correctness (24), security (16), performance (10), edge cases (8), resource (7), other (6). Difficulty: easy (18), medium (42), hard (31), expert (8). See Table~\ref{tab:dataset}.

\begin{table}[t]
\centering
\caption{Evaluation dataset composition (99 samples with verified labels; 70 Claude-generated, 29 hand-written).}
\label{tab:dataset}
\small
\begin{tabular}{@{}lrr@{}}
\toprule
\textbf{Category} & \textbf{Count} & \textbf{Percentage} \\
\midrule
\multicolumn{3}{l}{\textit{By Label}} \\
Buggy code (should FAIL) & 71 & 71.7\% \\
Correct code (should PASS) & 28 & 28.3\% \\
\midrule
\multicolumn{3}{l}{\textit{By Source}} \\
Hand-curated mirror & 29 & 29.3\% \\
Claude-generated & 70 & 70.7\% \\
\midrule
\multicolumn{3}{l}{\textit{By Bug Category (buggy samples)}} \\
Correctness bugs & 24 & 33.8\% \\
Security vulnerabilities & 16 & 22.5\% \\
Performance issues & 10 & 14.1\% \\
Edge case failures & 8 & 11.3\% \\
Resource management & 7 & 9.9\% \\
Other categories & 6 & 8.5\% \\
\midrule
\multicolumn{3}{l}{\textit{By Difficulty}} \\
Easy & 18 & 18.2\% \\
Medium & 42 & 42.4\% \\
Hard & 31 & 31.3\% \\
Expert & 8 & 8.1\% \\
\bottomrule
\end{tabular}
\end{table}

HumanEval~\cite{chen2021humaneval} tests functional correctness but lacks bug labels. SWE-bench (2,294)~\cite{jimenez2024swebench} has 29.6\% label errors~\cite{xia2025swebench}. We trade size for quality (100\% verified labels).

\subsection{Methodology}

Metrics: standard classification (accuracy, TPR, FPR, precision, F1). Confidence intervals via 1,000-iteration bootstrap~\cite{efron1994bootstrap}.

Paired baseline: Semgrep run with four rule packs (\texttt{p/security-audit}, \texttt{p/correctness}, \texttt{p/python}, \texttt{p/python-best-practices}) on the same hand-curated mirror samples. This is the one baseline evaluated on identical inputs, so it is the only direct head-to-head. For context we also cite published numbers for other methods (Meta Prompt Testing: 75\% TPR, 8.6\% FPR, test-based~\cite{wang2024metamorphic}; SAST benchmarks: 60--70\% detection, 30--40\% FPR~\cite{johnson2024sast, owasp2024}), but these use different datasets and methodologies and are not directly comparable.

\subsection{Ablation Design}

We test all 15 combinations: single agents (4), pairs (6), triples (4), full system (1). Each configuration is tested on the 29 hand-curated mirror samples (15 × 29 = 435 evaluations). The main evaluation (Section~6.1) uses all 99 samples (29 mirror + 70 Claude-generated). The ablation is descriptive: it attributes realized accuracy to each agent rather than testing a theoretical prediction. We report each agent's marginal contribution $\Delta_i = \mathbb{E}[\text{Acc}(\text{with } A_i)] - \mathbb{E}[\text{Acc}(\text{without } A_i)]$ averaged over subsets.

We also ran the system on 300 patches generated by Claude Sonnet 4.5 for SWE-bench Lite; these have no ground truth, so we report only the verdict distribution and latency. System: Python 3.10, asyncio. Code released: \url{https://github.com/ShreshthRajan/codex-verify}.

% ========================================================================
% SECTION 6: RESULTS
% ========================================================================

\section{Results}

We present the main evaluation against the paired Semgrep baseline, the ablation attributing detection to each agent, and deployment behavior on Claude Sonnet 4.5-generated patches.

\subsection{Main Evaluation Results}

Table~\ref{tab:main_results} presents our system's performance on the 99-sample benchmark compared to baselines.

\begin{table*}[t]
\centering
\caption{Main evaluation on 99 samples with verified labels. Semgrep is the paired baseline, run on the shared hand-curated mirror samples (28 samples) with four rule packs. Confidence interval via 1,000-iteration bootstrap. Meta Prompt and SAST-benchmark numbers are from the literature on other datasets and are not directly comparable.}
\label{tab:main_results}
\small
\begin{tabular}{@{}lcccc@{}}
\toprule
\textbf{System} & \textbf{Accuracy} & \textbf{TPR} & \textbf{FPR} & \textbf{F1 Score} \\
\midrule
Semgrep (4 rule packs, paired) & 28.6\% & 4.8\% & 0.0\% & 0.091 \\
\midrule
Meta Prompt Testing$^{\dagger}$ & --- & 75.0\% & 8.6\% & --- \\
SAST benchmarks$^{\dagger}$ & 60--70\% & 60--70\% & 30--40\% & --- \\
\midrule
\textsc{CodeX-Verify} (ours) & \textbf{68.7\%} $\pm$ 9.1\% & \textbf{76.1\%} & 50.0\% & \textbf{0.777} \\
\bottomrule
\end{tabular}
\vspace{1mm}
\footnotesize{$^{\dagger}$Literature numbers on different datasets/methodologies; shown for context only, not a paired comparison.}
\end{table*}

\textbf{Overall Performance.} \textsc{CodeX-Verify} reaches 68.7\% accuracy (95\% CI: [59.6\%, 77.8\%]) and 76.1\% TPR. On the shared hand-curated samples, Semgrep with four rule packs reaches 4.8\% TPR (1 of 21 bugs) and 28.6\% accuracy at 0\% FPR: it is precision-tuned and stays silent on logic, edge-case, and resource bugs. Coordinating four static agents recovers most of the bugs the single analyzer misses, at the cost of a far higher false-positive rate (50\% vs. 0\%). The system's TPR is close to the published Meta Prompt Testing figure (75\%), but that method runs tests on a different dataset and is not directly comparable.

\textbf{Confusion Matrix.} TP=54 (caught 54/71 bugs), TN=14 (accepted 14/28 good code), FP=14 (flagged 14/28 good code), FN=17 (missed 17/71 bugs). Precision = 79.4\% (when we flag code, it's buggy 79\% of the time), Recall = 76.1\% (we catch 76\% of bugs), F1 = 0.777.

\textbf{False Positives.} Our 50.0\% FPR (14/28) exceeds Meta Prompt's 8.6\% because we flag quality issues, not just functional bugs. Causes: 43\% missing exception handling (enterprise standard, not a functional bug), 29\% low edge case coverage (quality metric), 21\% flagging \texttt{import os} as dangerous (security conservatism), 7\% production readiness. These are design choices for enterprise deployment, not errors.

\textbf{By Category.} Table~\ref{tab:category}: 100\% detection on resource management (7/7), 87.5\% on security (7/8), 75\% on correctness (18/24), 60\% on performance (6/10), 0\% on edge cases (0/2, small sample).

\begin{table*}[t]
\centering
\caption{Performance by bug category on 99-sample evaluation.}
\label{tab:category}
\small
\begin{tabular}{@{}lrrr@{}}
\toprule
\textbf{Bug Category} & \textbf{Samples} & \textbf{Detected} & \textbf{Detection Rate} \\
\midrule
Resource management & 7 & 7 & 100.0\% \\
Async coordination & 1 & 1 & 100.0\% \\
Regex security & 1 & 1 & 100.0\% \\
State management & 1 & 1 & 100.0\% \\
\midrule
Security vulnerabilities & 8 & 7 & 87.5\% \\
Algorithmic complexity & 3 & 2 & 66.7\% \\
Correctness bugs & 24 & 18 & 75.0\% \\
Performance issues & 10 & 6 & 60.0\% \\
\midrule
Edge case logic & 2 & 0 & 0.0\% \\
Serialization security & 1 & 0 & 0.0\% \\
\bottomrule
\end{tabular}
\end{table*}

\subsection{Ablation Study}

Table~\ref{tab:ablation} shows results for all 15 agent combinations.

\begin{table*}[t]
\centering
\caption{Ablation study results across 15 configurations on 29 unique samples. Configurations ranked by accuracy. Agent abbreviations: C=Correctness, S=Security, P=Performance, St=Style.}
\label{tab:ablation}
\small
\begin{tabular}{@{}lcccc@{}}
\toprule
\textbf{Configuration} & \textbf{Agents} & \textbf{Accuracy} & \textbf{TPR} & \textbf{FPR} \\
\midrule
\multicolumn{5}{l}{\textit{Agent Pairs (n=2)}} \\
C + P & 2 & \textbf{79.3\%} & 83.3\% & 40.0\% \\
C + St & 2 & 75.9\% & 79.2\% & 40.0\% \\
C + S & 2 & 69.0\% & 70.8\% & 40.0\% \\
\midrule
\multicolumn{5}{l}{\textit{Single Agents (n=1)}} \\
Correctness & 1 & 75.9\% & 79.2\% & 40.0\% \\
Security & 1 & 20.7\% & 4.2\% & 0.0\% \\
Performance & 1 & 17.2\% & 0.0\% & 0.0\% \\
Style & 1 & 17.2\% & 0.0\% & 0.0\% \\
\midrule
\multicolumn{5}{l}{\textit{Agent Triples (n=3)}} \\
C + P + St & 3 & 79.3\% & 83.3\% & 40.0\% \\
C + S + P & 3 & 72.4\% & 75.0\% & 40.0\% \\
C + S + St & 3 & 69.0\% & 70.8\% & 40.0\% \\
S + P + St & 3 & 24.1\% & 8.3\% & 0.0\% \\
\midrule
\multicolumn{5}{l}{\textit{Full System (n=4)}} \\
C + S + P + St & 4 & 72.4\% & 75.0\% & 40.0\% \\
\midrule
\multicolumn{5}{l}{\textit{Other Pairs}} \\
S + P & 2 & 24.1\% & 8.3\% & 0.0\% \\
S + St & 2 & 20.7\% & 4.2\% & 0.0\% \\
P + St & 2 & 17.2\% & 0.0\% & 0.0\% \\
\bottomrule
\end{tabular}
\end{table*}

Average accuracy by agent count rises from 32.8\% (1 agent) to 47.7\% (2), 61.2\% (3), and 72.4\% (4). This averaging is dominated by the many weak subsets that exclude Correctness; the better reading is per-agent attribution below, not a single +39.7pp headline.

Correctness alone gets 75.9\% (strongest), while Security (20.7\%), Performance (17.2\%), and Style (17.2\%) are weak alone because each fires on a narrow bug class. Every subset without Correctness scores under 25\%, so Correctness provides the base coverage. The best pair (C+P: 79.3\%) and the best triple (C+P+St: 79.3\%) slightly exceed the full system (72.4\%): adding Security pulls accuracy down on this label mix because its conservative blocking raises false positives.

Marginal contributions, averaged over subsets: Correctness +53.9pp, Security $-5.2$pp, Performance $-1.5$pp, Style $-4.2$pp. The negative values are the honest finding: on aggregate accuracy, Security, Performance, and Style cost more in false positives than they recover in true positives. Their value is categorical, not aggregate: with all four agents the system catches 100\% of resource bugs (7/7) and 87.5\% of security bugs (7/8), which C+P alone does not. The deployment choice is therefore explicit: C+P for highest accuracy, all four agents for security and resource coverage.

\subsection{Positioning Against Baselines}

Figure~\ref{fig:comparison} places the system on the TPR-FPR plane. Only Semgrep is a paired comparison on identical samples; the Meta Prompt and SAST-benchmark points are literature figures on other datasets, shown for context.

\begin{figure*}[t]
\centering
\footnotesize
\begin{verbatim}
FPR
 |
50% +  CodeX-Verify (76% TPR, 50% FPR) <-- OURS
 |
35% +  SAST benchmarks (60-70% TPR, 30-40% FPR) [literature]
 |
10% +  Meta Prompt (75% TPR, 8.6% FPR) [test-based, literature]
 |
 0% +  Semgrep, 4 packs (4.8% TPR, 0% FPR) [paired, this work]
  +----+----+----+----+----+----+-> TPR
  0%  20%  40%  60%  75%  80% 100%
\end{verbatim}
\caption{TPR-FPR plane. Against Semgrep on identical samples (bottom), the system recovers most of the missed bugs at the cost of higher FPR. Dashed/literature points are not paired.}
\label{fig:comparison}
\end{figure*}

Against the paired Semgrep baseline, the system raises TPR from 4.8\% to 76.1\% and accuracy from 28.6\% to 68.7\%, while FPR rises from 0\% to 50\%. We do not run a significance test against the literature baselines, since they were not evaluated on these samples.

\subsection{Ablation Findings}

Figure~\ref{fig:ablation_scaling} shows scaling behavior.

\begin{figure*}[t]
\centering
\footnotesize
\begin{verbatim}
Accuracy
   |
80%+      /-- C+P (79.3%)
75%+     /  Full 4-agent (72.4%)
70%+    /
65%+   /
60%+  /  3-agent (61.2%) +13.5pp
55%+ /
50%+/  2-agent (47.7%) +14.9pp
45%+
40%+
35%+  1-agent (32.8%) +11.2pp
30%+
   +----+----+----+----+
       1    2    3    4  (agents)
Diminishing returns: +14.9pp, +13.5pp, +11.2pp
\end{verbatim}
\caption{Average accuracy by agent count. The curve is dominated by the many weak subsets that exclude Correctness; see per-agent attribution in the text.}
\label{fig:ablation_scaling}
\end{figure*}

\textbf{Finding 1: Correctness carries the detection.} Correctness alone reaches 75.9\% accuracy; no subset that excludes it exceeds 25\%. The system is, at its core, a correctness checker with three specialized add-ons.

\textbf{Finding 2: The add-on agents are categorical, not general.} Security, Performance, and Style have negative average marginal contributions on aggregate accuracy ($-5.2$, $-1.5$, $-4.2$pp): on general bugs they add more false positives than true positives. Their value is confined to their category. Security detects 87.5\% of security-specific bugs while contributing almost nothing to overall TPR, and the full system catches 100\% of resource bugs that C+P misses. The full system does not optimize accuracy; it buys category coverage at an accuracy cost.

\textbf{Finding 3: The best configuration depends on the goal.} C+P (79.3\% accuracy, 83.3\% TPR, 40\% FPR) is the most accurate configuration and beats the full system (72.4\%). For deployments that do not need security or resource coverage, C+P is the better choice; the full four-agent system is justified only when those categories matter.

\subsection{Real-World Validation on Claude Patches}

Table~\ref{tab:claude} reports system behavior on 300 Claude Sonnet 4.5-generated patches for SWE-bench Lite issues (no ground truth available).

\begin{table*}[t]
\centering
\caption{Verification verdicts on 300 Claude Sonnet 4.5-generated patches for SWE-bench Lite issues. No ground truth available (would require test execution); table reports system behavior distribution.}
\label{tab:claude}
\small
\begin{tabular}{@{}lrr@{}}
\toprule
\textbf{Verdict} & \textbf{Count} & \textbf{Percentage} \\
\midrule
PASS & 2 & 0.7\% \\
FAIL & 297 & 99.0\% \\
ERROR (parse failure) & 1 & 0.3\% \\
\bottomrule
\end{tabular}
\end{table*}

On 300 Claude patches the system flags 99.0\% as FAIL and passes 0.7\%. This near-universal FAIL rate is a warning sign, not a strength: with no ground truth we cannot tell how many are false positives, and the 50\% FPR measured on labeled data suggests many are. The result confirms the system's conservatism on realistic multi-line patches (median length far above the labeled samples) and should be read as a latency and behavior check, not a detection result. Mean latency is 17ms per patch.

% ========================================================================
% SECTION 7: DISCUSSION
% ========================================================================

\section{Discussion}

\subsection{What the Agents Contribute}

The agents cover different bug classes by design: Correctness finds logic errors and edge cases (75.9\% solo), Security finds injection and secrets (87.5\% on security bugs, 20.7\% overall), Performance finds complexity and resource issues (17.2\% overall), Style finds maintainability problems. They cover each other's blind spots: Correctness misses SQL injection, which Security catches; Security misses off-by-one errors, which Correctness catches. We assume but do not measure that these agents are close to conditionally independent given bug status.

Correctness alone gets 75.9\% while the full system gets 72.4\%. The drop is a real cost: adding Security, Performance, and Style makes the decision more conservative (Algorithm~\ref{alg:decision} blocks on a security HIGH), which raises false positives faster than it raises true positives on this label mix. The payoff is category coverage, not aggregate accuracy. The best pair (C+P, 79.3\%) beats both the strongest single agent and the full system; use C+P for accuracy, all four agents for security and resource coverage.

\subsection{False Positives}

Our 50\% FPR is the main limitation. Analyzing the 14 false positives: 43\% from missing exception handling (enterprise standard, not a functional bug), 29\% from low edge case coverage (quality metric), 21\% from flagging \texttt{import os} as dangerous (security conservatism), 7\% from production readiness. These are design choices for enterprise deployment: requiring exception handling prevents crashes; demanding edge case coverage reduces failures. Code lacking these may still work, explaining higher FPR than functional-only verification (Meta Prompt: 8.6\%).

Static analysis flags \emph{potential} issues (``might fail without exception handling'') while test execution checks \emph{actual} behavior (``did produce wrong output''). We flag security holes (SQL injection, secrets) and quality issues (missing docs) that tests miss. This trades higher FPR for detecting more bug types. Security-focused orgs use our strict mode; low-false-alarm orgs use test-based methods.

We tried reducing FPR: initial 80\% dropped to 50\% by downgrading style issues from MEDIUM to LOW. Further relaxation (allow 2+ security HIGH) cut FPR to 20\% but dropped TPR to 42\%. Our 76\% TPR, 50\% FPR is Pareto-optimal for static analysis; achieving 8.6\% FPR needs test execution.

The 50\% FPR works for enterprise security (finance, healthcare, infrastructure) where false alarms beat missed bugs. AWS Lambda gates, Google security review, and Microsoft SDL operate similarly. This limits use in permissive workflows where developer friction from false alarms outweighs benefits.

\subsection{Limitations}

\textbf{Sample Size.} $n=99$ gives ±9.1\% confidence intervals, wider than ideal. The labels are verified, which enables precise TPR/FPR measurement that noisy large benchmarks do not (SWE-bench has 29.6\% label errors~\cite{xia2025swebench}), but 99 samples is small and the ablation's per-agent attribution carries that uncertainty. The ablation used the 29-sample mirror subset; running it on all 99 would tighten it.

\textbf{Static Analysis.} We miss: (1) Dynamic bugs (race conditions, timing failures, state issues needing execution). (2) Wrong algorithms with correct structure (wrong logic but proper exception handling passes our checks). (3) Subtle semantic bugs (metamorphic testing~\cite{wang2024metamorphic} detects output inconsistencies we miss). These are fundamental static analysis limits, not implementation bugs. Hybrid static + test execution could fix this.

\textbf{Python Only.} We use Python AST and Python patterns (\texttt{pickle.loads}, Django SQL). Generalizing to C/C++, Java, TypeScript needs: (1) language AST parsers (tree-sitter supports 50+ languages), (2) pattern libraries (buffer overflows for C, type confusion for TypeScript), (3) re-calibration. The architecture and theory generalize, but agent internals need language-specific work.

\textbf{Dataset construction.} 70 of the 99 samples were generated by Claude and 29 were hand-written to mirror SWE-bench failure patterns; none are drawn directly from production LLM output with independent labels. The samples isolate bug patterns and are short (50--1500 characters, median 500) versus production files (100--1000 lines), and the 71\% buggy ratio is higher than observed rates (40--60\%~\cite{jimenez2024swebench}). The system was tuned on this set, so the reported accuracy is on-distribution. These are the main threats to external validity: the numbers should be read as evidence that agent coordination helps on controlled bug patterns, not as a production detection rate. The 300 Claude patches (Section~6.4) are more realistic but have no labels.

\subsection{Deployment Implications}

Layered verification: run \textsc{CodeX-Verify} (<200ms) as a fast triage gate, then a test-based method on what passes for functional correctness, with human review for borderline WARNING cases. This spends static-analysis latency (17ms/patch) before expensive test execution (2--5s/sample).

Security-critical: the system detects 87.5\% of security-specific bugs, which suits a pre-commit gate blocking security HIGH issues (Algorithm~\ref{alg:decision}) where a missed vulnerability costs more than a false alarm.

Developer-facing: the 50\% FPR causes alert fatigue. The C+P configuration (79.3\% accuracy, 40\% FPR) is a better default here, with WARNING shown as a non-blocking suggestion. The ablation gives the data to make this choice explicitly.

\subsection{Future Work}

\textbf{Hybrid Verification.} Combine static (\textsc{CodeX-Verify}, 200ms) with test-based (Meta Prompt, 5s): static triages, tests validate passing samples. Expected: 80--85\% TPR, 15--20\% FPR. Needs sandboxing and test generation.

\textbf{Learned Thresholds.} Our hand-tuned thresholds get 76\% TPR, 50\% FPR. Learning on 500+ samples via logistic regression, reinforcement learning, or multi-objective optimization could cut FPR by 10--15pp.

\textbf{Operating Point Analysis.} Our evaluation reports a single operating point (76\% TPR, 50\% FPR). Characterizing the full TPR-FPR tradeoff curve via threshold sweeps, ROC/PR curve analysis, and per-category calibration would enable deployment-specific operating point selection and cost-sensitive decision making.

\textbf{Multi-Language.} Adapting to C/C++, Java, TypeScript needs: (1) AST parsers (tree-sitter supports 50+ languages), (2) pattern libraries (buffer overflows, type confusion), (3) re-calibration. Architecture generalizes; agent internals need 2--3 weeks per language.

\textbf{Active Learning.} $n=99$ is below ideal $n \geq 127$. Active learning: train on 30 samples, query high-uncertainty cases, refine. Could hit ±7\% CI with $n \approx 70$ vs. $n \approx 150$ random, cutting labeling 50\%.


\subsection{Impact}

Our 76\% TPR cuts buggy code acceptance from 40--60\% to 24--36\%, enabling safer deployment in: (1) Code review (Copilot, Cursor, Tabnine), (2) Bug fixing (SWE-agent, AutoCodeRover), (3) Enterprise CI/CD. Sub-200ms latency enables real-time use.

Risks: Over-reliance could reduce human review, missing novel bugs. The 50\% FPR causes alert fatigue without good UX. Orgs might think 76\% TPR means ``catches all bugs''---24\% false negative rate means human oversight essential.

We release open-source (6,122 lines) for transparency. Hand-tuned thresholds embed human judgments about risk, potentially biasing toward specific security models.

\subsection{Lessons}

Curating 99 samples with verified labels (vs. SWE-bench's 2,294 with 29.6\% errors) enabled precise measurement. Quality beats quantity: smaller high-quality benchmarks give more reliable insights than large noisy ones. We trade $\pm$9.1\% vs. $\pm$3\% CI to eliminate label noise.

The ablation over all 15 subsets showed what a single aggregate number hides: Correctness carries the accuracy, and the other three agents trade accuracy for category coverage. The honest reading is that this is a correctness checker with useful specialized add-ons, not a uniform multi-agent gain.

Attempts to cut FPR below 50\% without tests all failed. Static analysis has precision ceilings: can't distinguish quality concerns from bugs without running code. Hybrid static + dynamic is the frontier.

% ========================================================================
% SECTION 8: CONCLUSION
% ========================================================================

\section{Conclusion}

LLM-generated code hides bugs that pass basic tests: 29.6\% of SWE-bench ``solved'' patches are behaviorally wrong, and 62\% of BaxBench backends are vulnerable or incorrect. \textsc{CodeX-Verify} is a static multi-agent verifier that coordinates four agents (correctness, security, performance, style) in under 200ms per sample without test execution.

On 99 samples with verified labels the system reaches 76.1\% TPR and 68.7\% accuracy. On the same samples, Semgrep with four rule packs reaches 4.8\% TPR and 28.6\% accuracy, so coordinating agents recovers most of the bugs a single static analyzer misses. The cost is a 50\% false-positive rate: the system flags quality defects that do not change program output, and this is its central limitation.

An ablation over all 15 agent subsets attributes the detection: Correctness carries the accuracy (75.9\% solo), while Security, Performance, and Style trade a small accuracy loss for coverage of security and resource bugs (100\% of resource bugs, 87.5\% of security bugs with all four agents). The best pair (C+P) reaches 79.3\% and is the right default when security coverage is not needed.

The theory is a single monotonicity fact (adding an informative agent does not reduce the system's information about bug presence); conditional independence of agents is assumed, not measured. The results hold on a 99-sample set that is partly self-generated and tuned on, so they are evidence that agent coordination helps on controlled bug patterns, not a production detection rate.

Three directions follow: hybrid static-dynamic verification to cut the false-positive rate with test execution; multi-language support via tree-sitter; and learned thresholds on a larger, independently labeled dataset.

% ========================================================================
% ACKNOWLEDGMENTS
% ========================================================================

\section*{Acknowledgments}
\small
We thank the reviewers for their constructive feedback. Code and data: \url{https://github.com/ShreshthRajan/codex-verify}.

% ========================================================================
% REFERENCES
% ========================================================================

\bibliographystyle{plain}
\bibliography{references}

% ========================================================================
% APPENDIX
% ========================================================================

\appendix

\section{Appendix A: Ablation Study Details}

Table~\ref{tab:ablation_full} shows detailed metrics for all 15 configurations, including precision, recall, F1, and execution time per configuration.

\begin{table*}[h]
\centering
\caption{Detailed ablation results for all 15 configurations with timing.}
\label{tab:ablation_full}
\small
\begin{tabular}{@{}lccccccc@{}}
\toprule
\textbf{Config} & \textbf{n} & \textbf{Acc} & \textbf{TPR} & \textbf{FPR} & \textbf{Prec} & \textbf{F1} & \textbf{Time (ms)} \\
\midrule
C+P & 2 & 79.3 & 83.3 & 40.0 & 83.3 & 0.833 & 95 \\
C+P+St & 3 & 79.3 & 83.3 & 40.0 & 83.3 & 0.833 & 120 \\
C & 1 & 75.9 & 79.2 & 40.0 & 79.2 & 0.792 & 82 \\
C+St & 2 & 75.9 & 79.2 & 40.0 & 79.2 & 0.792 & 105 \\
C+S+P+St & 4 & 72.4 & 75.0 & 40.0 & 75.0 & 0.750 & 148 \\
C+S+P & 3 & 72.4 & 75.0 & 40.0 & 75.0 & 0.750 & 135 \\
C+S & 2 & 69.0 & 70.8 & 40.0 & 70.8 & 0.708 & 110 \\
C+S+St & 3 & 69.0 & 70.8 & 40.0 & 70.8 & 0.708 & 128 \\
S+P+St & 3 & 24.1 & 8.3 & 0.0 & 100.0 & 0.154 & 98 \\
S+P & 2 & 24.1 & 8.3 & 0.0 & 100.0 & 0.154 & 85 \\
S & 1 & 20.7 & 4.2 & 0.0 & 100.0 & 0.080 & 68 \\
S+St & 2 & 20.7 & 4.2 & 0.0 & 100.0 & 0.080 & 78 \\
P & 1 & 17.2 & 0.0 & 0.0 & --- & 0.0 & 52 \\
St & 1 & 17.2 & 0.0 & 0.0 & --- & 0.0 & 58 \\
P+St & 2 & 17.2 & 0.0 & 0.0 & --- & 0.0 & 72 \\
\midrule
\textit{By agent count} \\
1 agent & --- & 32.8 & 20.8 & 10.0 & --- & --- & 65 \\
2 agents & --- & 47.7 & 41.0 & 20.0 & --- & --- & 92 \\
3 agents & --- & 61.2 & 59.4 & 30.0 & --- & --- & 120 \\
4 agents & --- & 72.4 & 75.0 & 40.0 & --- & --- & 148 \\
\bottomrule
\end{tabular}
\end{table*}

Configurations without Correctness achieve <25\% accuracy, demonstrating Correctness provides essential base coverage. Security/Performance/Style alone achieve 0\% TPR on general bugs but specialize in narrow domains (Security: 87.5\% detection on security-specific bugs). The best 2-agent pair (C+P) and best 3-agent configuration (C+P+St) achieve identical performance (79.3\%), indicating Style provides no marginal value when Correctness and Performance are present. Execution time scales sublinearly with agent count: 4 agents run in 148ms (parallel) vs. 260ms if run sequentially, achieving 1.76$\times$ speedup on average.

\section{Appendix B: Security Pattern Library}

Table~\ref{tab:security_patterns} lists the complete vulnerability detection pattern library used by the Security agent, with CWE mappings and base severity assignments.

\begin{table*}[h]
\centering
\caption{Vulnerability patterns with CWE mappings and severity.}
\label{tab:security_patterns}
\small
\begin{tabular}{@{}llll@{}}
\toprule
\textbf{Pattern} & \textbf{Example} & \textbf{CWE} & \textbf{Severity} \\
\midrule
SQL injection & \texttt{execute(...\%...)}, \texttt{f"SELECT \{x\}"} & CWE-89 & HIGH \\
Command injection & \texttt{os.system}, \texttt{shell=True} & CWE-78 & HIGH \\
Code execution & \texttt{eval()}, \texttt{exec()} & CWE-94 & CRITICAL \\
Unsafe deserialization & \texttt{pickle.loads()}, \texttt{yaml.load()} & CWE-502 & HIGH \\
Weak crypto & \texttt{md5()}, \texttt{sha1()}, \texttt{random.randint()} & CWE-327/338 & MEDIUM \\
Hardcoded secrets & \texttt{password = "..."}, \texttt{api\_key = "..."} & CWE-798 & HIGH \\
\bottomrule
\end{tabular}
\end{table*}

Context-aware severity escalation: SQL injection patterns near authentication keywords (\texttt{password}, \texttt{auth}, \texttt{login}) escalate from HIGH to CRITICAL with multiplier 2.5. Secret detection combines regex patterns (AWS keys, GitHub tokens, API keys) with entropy-based analysis ($H(s) > 3.5$ threshold for strings with length $|s| \geq 20$).

\section{Appendix C: Performance Characteristics}

Table~\ref{tab:timing} shows per-agent execution latency measurements across 99 samples, demonstrating the benefits of parallel execution.

\begin{table*}[h]
\centering
\caption{Per-agent execution latency breakdown showing parallelization benefits.}
\label{tab:timing}
\small
\begin{tabular}{@{}lrrr@{}}
\toprule
\textbf{Agent} & \textbf{Mean (ms)} & \textbf{Std (ms)} & \textbf{Max (ms)} \\
\midrule
Correctness & 82 & 18 & 150 \\
Security & 68 & 12 & 120 \\
Performance & 52 & 10 & 95 \\
Style & 58 & 8 & 88 \\
\midrule
Parallel (max) & 148 & 22 & 180 \\
Sequential (sum) & 260 & --- & 453 \\
\midrule
Speedup & 1.76$\times$ & --- & 2.52$\times$ \\
\bottomrule
\end{tabular}
\end{table*}

Parallel execution via \texttt{asyncio.gather()} achieves 1.76$\times$ average speedup over sequential execution (2.52$\times$ best case), with total latency bounded by the slowest agent (Correctness, 82ms mean). The sublinear scaling (4 agents in 148ms vs. 260ms sequential) validates the asynchronous architecture design.

\end{document}
